%% ---------------------------------------------------------------------
\section{Pasividad sin memoria}
\label{sec:pa-sinmemoria}
%% ---------------------------------------------------------------------

El caso más sencillo es un elemento sin estado, $y=h(u)$, como un resistor no lineal con $u$ el voltaje y $y$ la corriente. Sin estado no hay dónde almacenar ($V\equiv0$) y la desigualdad de disipación con $s=u\T y$ se reduce a que la potencia nunca sea negativa.

\begin{definicion}[label=def:pa-sinmemoria]{Pasividad sin memoria \citep[Def.~6.1]{Kha02}}
El sistema $y=h(u)$, $h:\R^m\to\R^m$, es
\begin{itemize}
  \item \clave{pasivo} si $u\T y\ge0$ para todo $u$;
  \item \clave{sin pérdidas} si $u\T y=0$ para todo $u$;
  \item \clave{pasivo con prealimentación de la entrada} si $u\T y\ge u\T\varphi(u)$ para alguna función $\varphi$, y \clave{estrictamente pasivo en la entrada} (ISP, \emph{input strictly passive}) si además $u\T\varphi(u)>0$ para todo $u\neq0$;
  \item \clave{pasivo con retroalimentación de la salida} si $u\T y\ge y\T\rho(y)$ para alguna función $\rho$, y \clave{estrictamente pasivo en la salida} (OSP, \emph{output strictly passive}) si además $y\T\rho(y)>0$ para todo $y\neq0$.
\end{itemize}
\end{definicion}

En el caso escalar, $u\,h(u)\ge0$ quiere decir que la gráfica de $h$ vive en el primer y tercer cuadrantes. El ``exceso'' $\varphi$ o $\rho$ dice cuánta potencia sobra; si $u\T\varphi(u)<0$ para algún $u$, hay \emph{escasez}, y el sistema puede no ser pasivo. Las dos definiciones con $\varphi$ y $\rho$ tienen una lectura de diagrama de bloques (\cref{fig:pa-transformaciones}): una \emph{prealimentación} $-\varphi(u)$ quita el exceso o la escasez en la entrada, y una \emph{retroalimentación positiva} $u=\tilde u+\rho(y)$ quita el de la salida. En ambos casos lo que queda es pasivo:
\[
  u\T\bigl(y-\varphi(u)\bigr)=u\T y-u\T\varphi(u)\ge0,
  \qquad
  \tilde u\T y=u\T y-\rho(y)\T y\ge0 .
\]

\begin{figure}[ht]
\centering
\includegraphics[width=\linewidth]{cap11/t_transformaciones}
\caption{Dos maneras de mover el exceso de pasividad. (a) Con $\varphi(u)=\epsilon u$, el sistema es pasivo si la gráfica de $h$ queda por encima de $\epsilon u$ en el primer cuadrante y por debajo en el tercero. (b) Con $\rho(y)=\delta y$, ídem para la gráfica inversa.}
\label{fig:pa-transformaciones}
\end{figure}

\begin{definicion}[label=def:pa-sector]{Sector}
Una función escalar $h$ está en el \clave{sector} $[\alpha,\beta]$, y se escribe $h\in[\alpha,\beta]$, si
\begin{equation}
  \alpha u^2\le u\,h(u)\le\beta u^2\qquad\text{para todo }u,
  \qquad\text{o, equivalentemente,}\qquad \bigl(h(u)-\alpha u\bigr)\bigl(h(u)-\beta u\bigr)\le0 .
  \label{eq:pa-sector}
\end{equation}
Con $\beta=\infty$ se escribe $h\in[\alpha,\infty]$ y se pide solo la primera desigualdad.
\end{definicion}

Si $h\in[0,\infty]$, $h$ es pasiva. Si $h\in[\epsilon,\infty]$ con $\epsilon>0$, es ISP con $\varphi(u)=\epsilon u$. Y si $h\in[0,\beta]$ con $\beta<\infty$, es OSP: como $h(u)$ tiene el signo de $u$ y $\abs{h(u)}\le\beta\abs u$,
\begin{equation}
  h(u)^2\;\le\;\beta\,\abs u\,\abs{h(u)}=\beta\,u\,h(u)
  \quad\Longrightarrow\quad
  u\,y\;\ge\;\tfrac1\beta\,y^2,\qquad \rho(y)=y/\beta .
  \label{eq:pa-sectorOSP}
\end{equation}
La saturación $\sat(u)$ está en el sector $[0,1]$: es pasiva \emph{y} OSP, con $u\sat(u)\ge\sat(u)^2$. Es la no linealidad del problema de Lur'e.

\begin{figure}[ht]
\centering
\includegraphics[width=\linewidth]{cap11/f_sin_memoria}
\caption{Tres no linealidades estáticas. Izquierda: la característica del diodo túnel de \citet[ej.~1.2.2]{Kha02}; tiene pendiente negativa y es pasiva. Centro: $h(u)=0.9u+0.25\sin2u$ está en el sector $[0.79,\,1.40]$, dentro de $[\alpha,\beta]=[0.25,\,1.5]$ (en verde). Derecha: la saturación, en el sector $[0,1]$.}
\label{fig:pa-sinmemoria}
\end{figure}

\begin{trampa}[label=tr:pa-pendiente]{``Pendiente negativa, luego no es pasivo''}
\plausible{El diodo túnel tiene resistencia diferencial negativa (pendiente negativa) entre $0.14$ y $0.66$\,V; por lo tanto, no es un elemento pasivo.}
La pasividad mira el \emph{signo de la potencia}, $u\,h(u)$, no la pendiente de $h$. La característica del diodo túnel tiene pendiente negativa en $(0.1381,\,0.6601)$\,V, pero su mínimo local vale $h(0.6601)=0.0683$\,mA$\,>0$: la gráfica nunca sale del primer y tercer cuadrantes y $u\,h(u)\ge0$ para todo $u$. El diodo no entrega energía; lo que hace la pendiente negativa es que, \emph{alrededor de un punto de operación}, las perturbaciones pequeñas vean una resistencia negativa. Por eso el circuito del capítulo~2 tiene un equilibrio inestable.
\correcto{La propiedad que mira las pendientes es otra: la \emph{pasividad incremental}, $(u_1-u_2)\bigl(h(u_1)-h(u_2)\bigr)\ge0$, que para $h\in\Cont1$ equivale a $h'\ge0$. El diodo túnel es pasivo y no es incrementalmente pasivo.}
\end{trampa}

\codigorepo{cap11_pasividad/C11_Ej3_SinMemoria.m}{diodo}{Pasivo con pendiente negativa}
